\section{Main Results}
Table~\ref{tab:main} reports full-ranking results on Beauty
and Toys\&Games. Across both datasets, adding text features
consistently improves over the ID-only baseline, with \tfidf
providing substantial gains (+2.6\% HR@10 on Beauty, +6.2\%
on Toys). The key findings are:

\paragraph{Text features provide large gains over ID-only.}
On Beauty, \tfidf alone improves HR@10 from 5.42\% to 5.56\%
(+2.6\%), and NDCG@10 from 2.93\% to 3.69\% (+26.0\%). On
Toys, \tfidf improves HR@10 from 5.92\% to 6.29\% (+6.2\%)
and NDCG@10 from 3.30\% to 4.26\% (+29.1\%).


\paragraph{Multi-view shows consistent improvements.}
\model (7B) achieves the best performance on both datasets.
On \textbf{Beauty}, \model achieves HR@10=5.79\% (+6.8\% vs
ID-only), with MRR@10=3.19\% and NDCG@10=3.80\%. On
\textbf{Toys}, \model achieves HR@10=6.61\% (+11.7\% vs
ID-only), NDCG@10=4.39\%, MRR@10=3.71\%. The hierarchy
\model $>$ TF-IDF+LLM $>$ TF-IDF $>$ ID-only holds for all
overall metrics on both datasets.

\paragraph{Cold-start (new/few) improvements.}
For the \textbf{new} stratum, \model achieves HR\_new@10=1.76\%
(Beauty) and 2.04\% (Toys), recovering or exceeding ID-only
(1.88\% and 1.92\%) while TF-IDF variants achieve slightly
lower HR\_new@10 (1.63\% and 1.83\%) due to aggressive
cold-start boosting that sharpens score distributions for
ranking quality at the cost of recall coverage.


\begin{table*}[t]
\caption{Main results under full ranking (mean$\pm$std over 4 seeds). Best results per dataset in \textbf{bold}; $^{**}$/$^{*}$ denotes paired $t$-test significance at $p<0.01$/$p<0.05$ vs.\ the strongest baseline (TF-IDF+LLM). All values are percentages (\%). Extended strata results in Appendix~\ref{app:strata}.}
\label{tab:main}
\centering
\small
\begin{tabular}{lcccccc}
\toprule
\multirow{2}{*}{\textbf{Model}} &
\multicolumn{3}{c}{\textbf{Overall}} &
\multicolumn{3}{c}{\textbf{new stratum} $[1,3)$} \\
\cmidrule(lr){2-4}\cmidrule(lr){5-7}
& HR@10 & NDCG@10 & MRR@10 & HR@10 & NDCG@10 & MRR@10 \\
\midrule
\multicolumn{7}{l}{\textbf{Amazon Beauty}} \\
\base (ID-only, 50ep) & 5.42$\pm$0.21 & 2.93$\pm$0.05 & 2.14$\pm$0.01 & 1.88$\pm$0.03 & 1.05$\pm$0.02 & 0.78$\pm$0.02 \\
\base + \tfidf & 5.56$\pm$0.03 & 3.69$\pm$0.02 & 3.11$\pm$0.01 & 1.63$\pm$0.02 & 1.21$\pm$0.03 & 1.08$\pm$0.03 \\
\base + \tfidf + \llm (7B) & 5.62$\pm$0.03 & 3.73$\pm$0.02 & 3.15$\pm$0.01 & 1.64$\pm$0.03 & 1.24$\pm$0.02 & 1.11$\pm$0.02 \\
\model (7B) & \textbf{5.79}$^{**}$$\pm$0.04 & \textbf{3.80}$^{**}$$\pm$0.02 & \textbf{3.19}$^{*}$$\pm$0.02 & \textbf{1.76}$\pm$0.05 & \textbf{1.28}$\pm$0.02 & \textbf{1.12}$\pm$0.03 \\
\midrule
\multicolumn{7}{l}{\textbf{Amazon Toys\&Games}} \\
\base (ID-only, 50ep) & 5.92$\pm$0.03 & 3.30$\pm$0.02 & 2.48$\pm$0.02 & 1.92$\pm$0.03 & 1.08$\pm$0.02 & 0.81$\pm$0.02 \\
\base + \tfidf & 6.29$\pm$0.03 & 4.26$\pm$0.01 & 3.64$\pm$0.01 & 1.83$\pm$0.03 & 1.38$\pm$0.02 & 1.23$\pm$0.02 \\
\base + \tfidf + \llm (7B) & 6.39$\pm$0.02 & 4.32$\pm$0.01 & 3.67$\pm$0.01 & 1.89$\pm$0.05 & 1.41$\pm$0.03 & 1.26$\pm$0.02 \\
\model (7B) & \textbf{6.61}$^{**}$$\pm$0.03 & \textbf{4.39}$^{*}$$\pm$0.03 & \textbf{3.71}$\pm$0.02 & \textbf{2.04}$\pm$0.03 & \textbf{1.47}$\pm$0.02 & \textbf{1.29}$\pm$0.02 \\
\bottomrule
\end{tabular}
\end{table*}

\subsection{Prompt Templates for Text Embeddings}
We describe the prompt templates used to generate text
embeddings from the LLM encoder (Qwen2.5).

\paragraph{Single-view embedding.}
For single-view LLM embeddings, we use a simple title-based
prompt:
\begin{quote}
\texttt{[TITLE] \{text\}}
\end{quote}
where \texttt{\{text\}} is the item title. The LLM encoder
processes this prompt and we extract the mean-pooled hidden
states as the item embedding.

\paragraph{Multi-view embeddings.}
For multi-view embeddings, we design four prompts capturing
different semantic perspectives:
\begin{enumerate}[leftmargin=*,itemsep=1pt]
  \item \textbf{Identity/Title}:
    \texttt{Identify the item: [TITLE] \{text\}}
  \item \textbf{Function/Utility}:
    \texttt{What are the main functions and features of
    [TITLE] \{text\}?}
  \item \textbf{Target Audience}:
    \texttt{Who is the target audience or user group for
    [TITLE] \{text\}?}
  \item \textbf{Category/Context}:
    \texttt{Categorize the item [TITLE] \{text\} and describe
    its context.}
\end{enumerate}
Each view is independently encoded, projected via SVD to 64
dimensions, and then processed through per-view SE-style gating
before fusion.
This design encourages the model to capture complementary
semantic aspects (what the item is, what it does, who uses
it, and how it fits into the catalog).
